
\documentclass[letterpaper,hide notes,xcolor={table,svgnames},pdftex,10pt]{beamer}
\def\showexamples{t}


%\usepackage[svgnames]{xcolor}

%% Demo talk
%\documentclass[letterpaper,notes=show]{beamer}

\usecolortheme{crane}
\setbeamertemplate{navigation symbols}{}

\usetheme{MyPittsburgh}
%\usetheme{Frankfurt}

%\usepackage{tipa}

\usepackage{hyperref}
\usepackage{graphicx,xspace}
\usepackage[normalem]{ulem}
\usepackage{multicol}
\usepackage{amsmath,amssymb,amsthm,graphicx,xspace}
\newcommand\SF[1]{$\bigstar$\footnote{SF: #1}}

\usepackage[default]{sourcesanspro}
\usepackage[T1]{fontenc}
\usepackage[scaled]{beramono}
\usepackage{tikzpagenodes}

\newcounter{tmpnumSlide}
\newcounter{tmpnumNote}


% old question code
%\newcommand\question[1]{{$\bigstar$ \small \onlySlide{2}{#1}}}
% \newcommand\nquestion[1]{\ifdefined \presentationonly \textcircled{?} \fi \note{\par{\Large \textbf{?}} #1}}
% \newcommand\nanswer[1]{\note{\par{\Large \textbf{A}} #1}}


 \newcommand\mnote[1]{%
   \addtocounter{tmpnumSlide}{1}
   \ifdefined\showcues {~\tiny\fbox{\arabic{tmpnumSlide}}}\fi
   \note{\setlength{\parskip}{1ex}\addtocounter{tmpnumNote}{1}\textbf{\Large \arabic{tmpnumNote}:} {#1\par}}}

\newcommand\mmnote[1]{\note{\setlength{\parskip}{1ex}#1\par}}

%\newcommand\mnote[2][]{\ifdefined\handoutwithnotes {~\tiny\fbox{#1}}\fi
% \note{\setlength{\parskip}{1ex}\textbf{\Large #1:} #2\par}}

%\newcommand\mnote[2][]{{\tiny\fbox{#1}} \note{\setlength{\parskip}{1ex}\textbf{\Large #1:} #2\par}}

\newcommand\mquestion[2]{{~\color{red}\fbox{?}}\note{\setlength{\parskip}{1ex}\par{\Large \textbf{?}} #1} \note{\setlength{\parskip}{1ex}\par{\Large \textbf{A}} #2\par}\ifdefined \presentationonly \pause \fi}

\newcommand\blackboard[1]{%
\ifdefined   \showblackboard
  {#1}
  \else {\begin{center} \fbox{\colorbox{blue!30}{%
         \begin{minipage}{.95\linewidth}%
           \hspace{\stretch{1}} Some space intentionally left blank; done at the blackboard.%
         \end{minipage}}}\end{center}}%
         \fi%
}



%\newcommand\q{\tikz \node[thick,color=black,shape=circle]{?};}
%\newcommand\q{\ifdefined \presentationonly \textcircled{?} \fi}

\usepackage{listings}
\lstset{basicstyle=\footnotesize\ttfamily,
	breaklines=true,
	aboveskip=15pt,
  	belowskip=15pt,
	frame=lines,
	numbers=left, basicstyle=\scriptsize, numberstyle=\tiny, stepnumber=0, numbersep=2pt
}

\usepackage{siunitx}
\newcommand\sius[1]{\num[group-separator = {,}]{#1}\si{\micro\second}}
\newcommand\sims[1]{\num[group-separator = {,}]{#1}\si{\milli\second}}
\newcommand\sins[1]{\num[group-separator = {,}]{#1}\si{\nano\second}}
\sisetup{group-separator = {,}, group-digits = true}

%% -------------------- tikz --------------------
\usepackage{tikz}
\usetikzlibrary{positioning}
\usetikzlibrary{arrows,backgrounds,automata,decorations.shapes,decorations.pathmorphing,decorations.markings,decorations.text,decorations.pathreplacing}

\tikzstyle{place}=[circle,draw=blue!50,fill=blue!20,thick, inner sep=0pt,minimum size=6mm]
\tikzstyle{transition}=[rectangle,draw=black!50,fill=black!20,thick, inner sep=0pt,minimum size=4mm]

\tikzstyle{block}=[rectangle,draw=black, thick, inner sep=5pt]
\tikzstyle{bullet}=[circle,draw=black, fill=black, thin, inner sep=2pt]

\tikzstyle{pre}=[<-,shorten <=1pt,>=stealth',semithick]
\tikzstyle{post}=[->,shorten >=1pt,>=stealth',semithick]
\tikzstyle{bi}=[<->,shorten >=1pt,shorten <=1pt, >=stealth',semithick]

\tikzstyle{mut}=[-,>=stealth',semithick]

\tikzstyle{treereset}=[dashed,->, shorten >=1pt,>=stealth',thin]

\usepackage{ifmtarg}
\usepackage{xifthen}
\makeatletter
% new counter to now which frame it is within the sequence
\newcounter{multiframecounter}
% initialize buffer for previously used frame title
\gdef\lastframetitle{\textit{undefined}}
% new environment for a multi-frame
\newenvironment{multiframe}[1][]{%
\ifthenelse{\isempty{#1}}{%
% if no frame title was set via optional parameter,
% only increase sequence counter by 1
\addtocounter{multiframecounter}{1}%
}{%
% new frame title has been provided, thus
% reset sequence counter to 1 and buffer frame title for later use
\setcounter{multiframecounter}{1}%
\gdef\lastframetitle{#1}%
}%
% start conventional frame environment and
% automatically set frame title followed by sequence counter
\begin{frame}%
\frametitle{\lastframetitle~{\normalfont(\arabic{multiframecounter})}}%
}{%
\end{frame}%
}
\makeatother

\makeatletter
\newdimen\tu@tmpa%
\newdimen\ydiffl%
\newdimen\xdiffl%
\newcommand\ydiff[2]{%
    \coordinate (tmpnamea) at (#1);%
    \coordinate (tmpnameb) at (#2);%
    \pgfextracty{\tu@tmpa}{\pgfpointanchor{tmpnamea}{center}}%
    \pgfextracty{\ydiffl}{\pgfpointanchor{tmpnameb}{center}}%
    \advance\ydiffl by -\tu@tmpa%
}
\newcommand\xdiff[2]{%
    \coordinate (tmpnamea) at (#1);%
    \coordinate (tmpnameb) at (#2);%
    \pgfextractx{\tu@tmpa}{\pgfpointanchor{tmpnamea}{center}}%
    \pgfextractx{\xdiffl}{\pgfpointanchor{tmpnameb}{center}}%
    \advance\xdiffl by -\tu@tmpa%
}
\makeatother
\newcommand{\copyrightbox}[3][r]{%
\begin{tikzpicture}%
\node[inner sep=0pt,minimum size=2em](ciimage){#2};
\usefont{OT1}{phv}{n}{n}\fontsize{4}{4}\selectfont
\ydiff{ciimage.south}{ciimage.north}
\xdiff{ciimage.west}{ciimage.east}
\ifthenelse{\equal{#1}{r}}{%
\node[inner sep=0pt,right=1ex of ciimage.south east,anchor=north west,rotate=90]%
{\raggedleft\color{black!50}\parbox{\the\ydiffl}{\raggedright{}#3}};%
}{%
\ifthenelse{\equal{#1}{l}}{%
\node[inner sep=0pt,right=1ex of ciimage.south west,anchor=south west,rotate=90]%
{\raggedleft\color{black!50}\parbox{\the\ydiffl}{\raggedright{}#3}};%
}{%
\node[inner sep=0pt,below=1ex of ciimage.south west,anchor=north west]%
{\raggedleft\color{black!50}\parbox{\the\xdiffl}{\raggedright{}#3}};%
}
}
\end{tikzpicture}
}


%% --------------------

%\usepackage[excludeor]{everyhook}
%\PushPreHook{par}{\setbox0=\lastbox\llap{MUH}}\box0}

%\vspace*{\stretch{1}

%\setbox0=\lastbox \llap{\textbullet\enskip}\box0}

\setlength{\parskip}{\fill}

\newcommand\noskips{\setlength{\parskip}{1ex}}
\newcommand\doskips{\setlength{\parskip}{\fill}}

\newcommand\xx{\par\vspace*{\stretch{1}}\par}
\newcommand\xxs{\par\vspace*{2ex}\par}
\newcommand\tuple[1]{\langle #1 \rangle}
\newcommand\code[1]{{\sf \footnotesize #1}}
\newcommand\ex[1]{\uline{Example:} \ifdefined \presentationonly \pause \fi
  \ifdefined\showexamples#1\xspace\else{\uline{\hspace*{2cm}}}\fi}

\newcommand\ceil[1]{\lceil #1 \rceil}


\AtBeginSection[]
{
   \begin{frame}
       \frametitle{Outline}
       \tableofcontents[currentsection]
   \end{frame}
}



\pgfdeclarelayer{edgelayer}
\pgfdeclarelayer{nodelayer}
\pgfsetlayers{edgelayer,nodelayer,main}

\tikzstyle{none}=[inner sep=0pt]
\tikzstyle{rn}=[circle,fill=Red,draw=Black,line width=0.8 pt]
\tikzstyle{gn}=[circle,fill=Lime,draw=Black,line width=0.8 pt]
\tikzstyle{yn}=[circle,fill=Yellow,draw=Black,line width=0.8 pt]
\tikzstyle{empty}=[circle,fill=White,draw=Black]
\tikzstyle{bw} = [rectangle, draw, fill=blue!20, 
    text width=4em, text centered, rounded corners, minimum height=2em]
    
    \newcommand{\CcNote}[1]{% longname
	This work is licensed under the \textit{Creative Commons #1 3.0 License}.%
}
\newcommand{\CcImageBy}[1]{%
	\includegraphics[scale=#1]{creative_commons/cc_by_30.pdf}%
}
\newcommand{\CcImageSa}[1]{%
	\includegraphics[scale=#1]{creative_commons/cc_sa_30.pdf}%
}
\newcommand{\CcImageNc}[1]{%
	\includegraphics[scale=#1]{creative_commons/cc_nc_30.pdf}%
}
\newcommand{\CcGroupBySa}[2]{% zoom, gap
	\CcImageBy{#1}\hspace*{#2}\CcImageNc{#1}\hspace*{#2}\CcImageSa{#1}%
}
\newcommand{\CcLongnameByNcSa}{Attribution-NonCommercial-ShareAlike}

\newenvironment{changemargin}[1]{% 
  \begin{list}{}{% 
    \setlength{\topsep}{0pt}% 
    \setlength{\leftmargin}{#1}% 
    \setlength{\rightmargin}{1em}
    \setlength{\listparindent}{\parindent}% 
    \setlength{\itemindent}{\parindent}% 
    \setlength{\parsep}{\parskip}% 
  }% 
  \item[]}{\end{list}} 




\title{Lecture 19 --- Single Thread Performance}

\author{Patrick Lam \\ \small \texttt{patrick.lam@uwaterloo.ca}}
\institute{Department of Electrical and Computer Engineering \\
  University of Waterloo}
\date{\today}


\begin{document}

\begin{frame}
  \titlepage

\end{frame}

\part{Single-Thread Performance}



\begin{frame}
\frametitle{Single-Thread Performance}

``Can you run faster just by trying harder?''

\begin{center}
\includegraphics[width=0.7\textwidth]{images/theflash.jpg}
\end{center}

\end{frame}


\begin{frame}
\frametitle{Single-Thread Performance}
\Large
\begin{changemargin}{2cm}
Performance improvements to date have \\
used parallelism
to improve throughput. 

Decreasing latency is trickier---\\
often requires domain-specific
tweaks. 

Today: one example of decreasing latency: \\
\hspace*{2em} Stream VByte.
\end{changemargin}

\end{frame}


\begin{frame}
\frametitle{I have a cunning plan...}

\Large
\begin{changemargin}{2cm}
Even Stream VByte uses parallelism:\\
\hspace*{2em}vector instructions. 

But there
are sequential improvements, \\
e.g. Stream VByte takes care to be predictable
for the branch predictor.
\end{changemargin}

\end{frame}
\begin{frame}
\frametitle{Inverted Indexes (like it's CS 137 again!)}

\Large
\begin{changemargin}{1cm}
\vspace*{-1em}
Abstractly: store a sequence of small integers.


Why Inverted indexes?

\hspace*{1cm}allow fast lookups by term;\\
\hspace*{1cm}support boolean queries combining terms.

\end{changemargin}

\end{frame}
\begin{frame}
\frametitle{Dogs, cats, cows, goats. In ur documents.}
\Large
\begin{changemargin}{2cm}
\begin{center}
\begin{tabular}{r|l}
docid & terms \\ \hline
1 & dog, cat, cow\\
2 & cat\\
3 & dog, goat\\
4 & cow, cat, goat\\
\end{tabular}
\end{center}
\end{changemargin}

\end{frame}


\begin{frame}
\frametitle{Inverting the Index}

\large\begin{changemargin}{2cm}
Here's the index and the inverted index:
\begin{center}
\begin{tabular}{r|l}
docid & terms \\ \hline
1 & dog, cat, cow\\
2 & cat\\
3 & dog, goat\\
4 & cow, cat, goat\\
\end{tabular} \hspace*{2em}
\begin{tabular}{r|l}
term & docs \\ \hline
dog & 1, 3 \\
cat & 1, 2, 4 \\
cow & 1, 4 \\
goat & 3, 4
\end{tabular}
\end{center}

Inverted indexes contain many small integers.

Deltas typically small if doc ids are sorted.
\end{changemargin}

\end{frame}


\begin{frame}
\frametitle{Storing inverted index lists: VByte}

\Large\begin{changemargin}{2cm}
\vspace*{-2em}
VByte uses a variable number of bytes\\
to store integers.  

Why? Most integers are
small,\\
especially on today's 64-bit processors.
\end{changemargin}
\end{frame}


\begin{frame}
\frametitle{How VByte Works}

\large\begin{changemargin}{1cm}
VByte works like this:

\begin{itemize}
\item $x$ between 0 and $2^7-1$ (e.g. $17 = 0b10001$):\\
\hspace*{1em}$0xxx xxxx$, e.g. $0001 0001$;
\item $x$ between $2^7$ and $2^{14}-1$ (e.g. $1729 = 0b110 11000001$):\\
\hspace*{1em}                   $1xxx xxxx/0xxx xxxx$ (e.g. $1100 0001/0000 1101$);\\
\item $x$ between $2^{14}$ and $2^{21}-1$: \\
\hspace*{1em}$0xxx xxxx/1xxx xxxx/1xxx xxxx$;
\item etc.
\end{itemize}

Control bit, or high-order bit, is:\\
\hspace*{2em}0 once done representing the int,\\
\hspace*{2em}1 if more bits remain.
\end{changemargin}

\end{frame}



\begin{frame}
\frametitle{Why VByte Helps}

\large\begin{changemargin}{2cm}
Isn't dealing with variable-byte integers harder?\\
\hspace*{2em} $\bullet$~ Yup!

But perf improves: \\
\hspace*{2em} $\bullet$~  We are using fewer bits! 

We fit more information into RAM and
cache, and can get higher throughput. (think inlining)

Storing and reading 0s isn't good use of resources. 

However, a naive algorithm to decode VByte \\
gives branch mispredicts.
\end{changemargin}
\end{frame}

\begin{frame}
\frametitle{Stream VByte}
\vspace*{-2em}
\large\begin{changemargin}{2cm}
Stream VByte: a variant of VByte using SIMD.


Science is incremental. \\
Stream VByte builds on earlier work---\\
\hspace*{2em}masked VByte, {\sc varint}-GB, {\sc varint}-G8IU. 

Innovation in Stream VByte:\\
\hspace*{2em}\emph{store the control and data streams separately}.
\end{changemargin}

\end{frame}

\begin{frame}
\frametitle{Control Stream}
\vspace*{-4em}
\large\begin{changemargin}{2cm}
Stream VByte's control stream uses two bits per integer to represent the size of the integer:
\begin{center}
\vspace*{-3em}
\begin{tabular}{ll@{~~~~~~~~}ll}
00 & 1 byte & 10 & 3 bytes\\
01 & 2 bytes & 11 & 4 bytes
\end{tabular}
\end{center}
\end{changemargin}

\end{frame}


\begin{frame}
\frametitle{Decoding Stream VByte}

\vspace*{-4em}
\large\begin{changemargin}{2cm}
Per decode iteration:\\
\hspace*{2em} reads 1 byte from the control stream,\\
\hspace*{2em} and 16 bytes of data.


Lookup table on control stream byte: decide how many
bytes it needs out of the 16 bytes it has read.

SIMD instructions: \\
\hspace*{2em}shuffle the bits each into their own integers. 

Unlike VByte, \\
Stream VByte uses all 8 bits
of data bytes as data.
\end{changemargin}

\end{frame}



\begin{frame}
\frametitle{Stream VByte Example}

\large\begin{changemargin}{2cm}
Say control stream contains $0b1000~1100$. \\
Then the data stream
contains the following sequence of integer sizes: $3, 1, 4, 1$. 

Out of the 16 bytes read,
this iteration uses 9 bytes; \\
\hspace*{2em} $\Rightarrow$ it advances the data pointer by 9. 

The SIMD
``shuffle'' instruction puts decoded integers from data stream at known positions in the
128-bit SIMD register.

Pad the first 3-byte integer with 1 byte, then
the next 1-byte integer with 3 bytes, etc. 
\end{changemargin}
\end{frame}

\begin{frame}
\frametitle{Stream VByte: Shuffling the Bits}
\vspace*{-1em}
\large\begin{changemargin}{1cm}
Say the data input is:\\
{\tt 0xf823~e127~2524~9748~1b..~....~....~....}. 

The 128-bit output is:\\
{\tt 0x00f8~23e1/0000~0027/2524 9748/0000/001b}\\
/s denote separation
between outputs. 

Shuffle mask is precomputed and
read from an array.
\end{changemargin}
\end{frame}

\begin{frame}[fragile]
\frametitle{SIMD Instructions}

\vspace*{-1em}
\large\begin{changemargin}{1cm}
The core of the implementation uses\\
three SIMD instructions:
\begin{lstlisting}[language=C]
  uint8_t C = lengthTable[control];
  __m128i Data = _mm_loadu_si128 ((__m128i *) databytes);
  __m128i Shuf = _mm_loadu_si128(shuffleTable[control]);
  Data = _mm_shuffle_epi8(Data, Shuf);
  databytes += C; control++;
\end{lstlisting}
\end{changemargin}
\end{frame}


\begin{frame}
\frametitle{But Does It Work?!}

\vspace*{-1em}
\large\begin{changemargin}{1cm}
Stream VByte performs better than previous techniques on a realistic input.


Why?

\begin{itemize}
\item control bytes are sequential:\\
\hspace*{1em} CPU can always prefetch the next control byte, \\
\hspace*{1em} because
its location is predictable;
\item data bytes are sequential \\
\hspace*{1em}and loaded at high throughput;
\item shuffling exploits the instruction set: \\
\hspace*{1em}takes 1 cycle;
\item control-flow is regular \\
\hspace*{1em}(tight loop which retrieves/decodes control
\& data;\\
\hspace*{1em}no conditional jumps).
\end{itemize}
\end{changemargin}
\end{frame}



\end{document}

